\section{Physical formulation and acceptance}
\label{app:physical}
For readability, subscripts $z,t$ are omitted where unambiguous. Let $d$ be
demand, $q^-$ deferred demand and $q^+$ recovered demand. With line incidence
$A_{z\ell}$, the balance is
\begin{multline}
 p^{\rm th}+p^{\rm nuc}+p^{\rm ren}+p^{\rm hydro}
 +p^{\rm bat,dis}+p^{\rm pump,dis}
 +\sum_\ell A_{z\ell}f_\ell+p^{\rm imp}+p^{\rm uns}\\
 =d-q^-+q^++p^{\rm bat,ch}+p^{\rm pump,ch}+p^{\rm exp}+p^{\rm excess}.
\end{multline}
Imports and exports enter the designated trade-anchor zone; a system mode
prevents simultaneous import/export. Each line has signed flow bounded by
its capacity. Renewable output is bounded by its realized availability;
unused availability is recorded as curtailment.

Thermal commitment $u$, start-up $v$ and shutdown $w$ satisfy
\begin{align}
 u_t-u_{t-1}&=v_t-w_t,&v_t+w_t&\leq1,\\
 \underline P u_t&\leq p_t\leq\overline P u_t,\\
 p_t-p_{t-1}&\leq R u_{t-1}+S^{\rm up}v_t,&
 p_{t-1}-p_t&\leq R u_t+S^{\rm down}w_t.
\end{align}
Initial output and commitment are fixed scenario data. Nuclear output obeys
its installed capacity and minimum operating limits. The model represents
aggregate zonal decisions, not an independently committed binary per graph asset.

For each battery or pumped store, energy satisfies
\begin{equation}
 e_{t+1}=\rho e_t+\Delta t\left(\eta^{\rm ch}p_t^{\rm ch}
                        -p_t^{\rm dis}/\eta^{\rm dis}\right),
\end{equation}
with power and energy bounds, mutually exclusive charging/discharging modes,
initial energy and specified terminal lower/upper bounds. Reservoir energy
additionally accounts for inflow, release and spill, with conversion
efficiency and reservoir/generation capacity bounds.
Demand response has causal backlog
\begin{equation}
 h_{t+1}=h_t+\Delta t(q_t^--q_t^+),\qquad h_t\geq0,\quad h_0=h_T=0.
\end{equation}
Deferral and recovery obey their power limits, activation restrictions and
ramping constraints. A terminal equality prevents deferred demand from
becoming unreported load shedding. Activation and block costs are included
in the physical objective.

Acceptance reconstructs these equations and costs from the dispatch, checks
finite values and axis consistency, and checks binary integrality separately.
ENS remains physically permitted and penalized. The $10^{-3}$ MWh service
cap is an aggregated outcome tolerance, distinct from per-constraint physical
and integrality tolerances.

\section{Graph representation and learned strategies}
\label{app:learned}
\subsection{Graph and encoder contract}
The graph contains system, region, zone, asset and weather nodes. Spatial
relations include hierarchical containment, weather influence and transmission;
temporal copies and relation templates describe storage, ramping and demand
flexibility. Node attributes summarize capacities, exogenous trajectories,
operating parameters and node type. The graph builder orders zones
lexicographically; embeddings are explicitly reordered to the physical
zone axis before decoding.

The encoder uses sparse graph-attention layers for spatial aggregation,
hierarchical pooling, a temporal Transformer and top-down context propagation.
The deployed encoder receives node features, connectivity, node types and
hierarchical assignments. Edge-type identifiers organize construction and
hierarchy; they are not a separate learned relation-specific input at this
interface. Transmission-distance metadata is not consumed by the encoder.
This distinction prevents attributing predictive information to every
attribute shown in a conceptual graph.

Input standardization uses frozen training moments, including a checked
categorical weather encoding. No evaluation label is used to construct an
embedding. Encoder pretraining membership, supervised strategy splits and
normalization provenance are separate audit objects: checking one does not
certify the others.

\subsection{Historical autoencoder, predictor and energy model}
A bidirectional GRU encodes each zone's strategy trajectory. Time pooling
and learned query-token cross-attention aggregate zones into 32 tokens of
dimension 192. The decoder uses learned zone/time queries, cross-attention
to the tokens and conditioning by 128-dimensional scenario embeddings. Its
capacity is 300 zones and 24 periods. These limits do not establish the
empirically observed training support.

The autoencoder minimizes weighted binary cross-entropy (BCE), with positive
class weights estimated on training data and capped at five. The predictor
maps scenario embeddings to strategy tokens, with the autoencoder fixed.
Its objective combines latent MSE, a variance term, decoded strategy BCE,
auxiliary decision BCE and temporal regularization, with respective weights
$1$, $0.1$, $5$, $2$ and $0.05$. The configured covariance penalty is zero.
This supervised strategy-prediction objective is distinct from image JEPA.
Training uses batch size eight, AdamW with initial learning rate $3\times10^{-4}$,
weight decay $10^{-4}$, gradient clipping at one, and two warm-up epochs.
AE/predictor epoch ceilings are 60/50 and early-stopping patience is 10/8.
Validation macro-F1 selects the checkpoints; the selected epoch indices are
29 and 36, respectively. The common supervised split has 4,500 training and
500 validation cases; the 600 evaluation scenarios are a separate cohort.

The energy model learns to rank reference latents below noisy or sampled
negative latents using a softplus contrastive objective with energy
regularization $10^{-3}$. It uses scenario context, hidden width 256 and
bounded energies. Training has the same learning rate and weight decay,
a 20-epoch ceiling, patience six, and validation energy-gap selection.
The selected epoch index is ten. None of these validation criteria directly
optimizes end-to-end service coverage.

Both historical learned methods apply channel thresholds
\begin{equation}
 (0.20,0.30,0.20,0.30,0.20,0.35,0.45,0.80)
\end{equation}
in the order of Eq.~\eqref{eq:strategy}. They were selected on validation
data from a 0.20--0.80 grid with 0.05 increments and frozen for evaluation.
The same threshold contract is retained by the no-gradient control.

For anchor $a$ and candidate latent $z_k$, a sampling step is
\begin{equation}
 z_{k+1}=\operatorname{clip}_{[-6,6]}\!\left[
 z_k-\epsilon_k\widetilde\nabla E(z_k,a,h)
 -0.02\epsilon_k(z_k-a)
 +\sqrt{2\epsilon_k}\,\sigma_k\xi_k\right],
 \quad \xi_k\sim\mathcal N(0,I).
\end{equation}
The gradient is elementwise clipped at one and then normalized by its
candidate-wise Euclidean norm with floor $10^{-6}$. Initial step and noise
scales are 0.02 and 0.05 with cosine annealing over twenty updates. Initialization
mixes anchor perturbations of standard deviation 0.25 with a latent prior,
using prior probability 0.5. Candidates have independent seeded random
streams. The no-gradient control sets only $\widetilde\nabla E$ to zero,
preserving the draws, drift, clipping and remaining operations.

\subsection{Repair and operational timing}
The first repair stage fixes the supplied seven-channel targets while
retaining structural fixes and the original domains of other variables.
The next stages release selected critical decisions, with a deviation penalty
of $10^4$ and configured flip limits 100 and 1,000. Up to three stages are
allowed. The cascade's internal slack trigger is 1 MWh; this internal trigger
does not replace physical validation or the stricter service endpoint.
Physical-cost selection compares only candidates validated and received
before the computation cutoff.

The finalization reserve was calibrated on 30 validation scenarios and four
methods, giving 114 observed natural finalizations and six deadline failures.
Twice the maximum observed finalization time, with a 0.05 s lower bound,
gives $r_{120}=38.4596436$ s. Failures with unobserved finalization are not
assigned zero time. For budget $b$, the unchanged rule is
\begin{equation}
 r(b)=\min\{0.5b,\max(0.05,r_{120}\min(1,b/120))\}.
\end{equation}
All online work must finish early enough to write and close normally before
$b$. Static workers contain no evaluation scenario. The supervisor measures
actual process-tree termination; forced termination never becomes a success.
Stage timings for concurrent candidates are reported alongside, not summed
into, elapsed wall-clock time.

\section{Controlled changes and the variable-horizon extension}
\label{app:extensions}
New development and final base identifiers use separate deterministic seed
namespaces. The audit checks UUIDs, realized physical inputs and shared base
groups. All horizons of a training base belong to the same split, with a
deterministic one-in-five validation rule. Labels are independently valid
incumbents; certification is recorded separately, and failed labeling attempts
are explicit exclusions, not silently redrawn cases.

Network recabling preserves the number of lines and their capacity/distance
multiset while ensuring connectivity. The fleet change multiplies thermal
capacity, minimum output, ramp and transition limits by 0.8 and battery power,
energy and initial/terminal bounds by 1.2. Duplication copies zones and
resources, doubles import capacity and connects the copies with a bridge.
Its starting bases have at most 150 zones. Permutation consistently remaps
all zone-indexed data, line endpoints and the import anchor.

For extended horizons, daily demand multipliers are 1, 0.95 and 1.05;
the same factors apply to demand-response limits. Solar, wind and other
renewable availability use daily multipliers 1, 0.90 and 0.85.
Other repeated profiles retain their first-day values.
Hourly resolution is unchanged, and terminal requirements apply at the new
endpoint. These controls should not be interpreted as realistic weather
forecast scenarios.

The new autoencoder uses a bidirectional GRU and token attention but replaces
learned zone/time-index queries with projections of scenario embeddings plus
sinusoidal time positions. Its latent predictor pools scenario sequences;
the energy network combines a candidate, its anchor and averaged scenario
context. Token count/dimension remain 32/192, with four attention heads.
The frozen graph encoder's sinusoidal buffer is extended analytically to
72 periods; its weights and normalization are unchanged.

Training uses single-instance batches with accumulation over eight examples,
AdamW at $2\times10^{-4}$ with weight decay $10^{-4}$, gradient clipping at one
and 60/50/20 epoch ceilings.
The AE minimizes BCE with positive weights capped at twenty; the predictor
minimizes latent MSE plus decoded BCE; the energy model minimizes
$\operatorname{softplus}(E^+-E^-)+10^{-3}((E^+)^2+(E^-)^2)$.
Validation loss selects each checkpoint. All channels use threshold 0.5.
The sampler uses twenty updates, step size 0.02, additive noise scale 0.05,
initial perturbation 0.15, elementwise gradient clipping at one and latent
clipping at ten. Its matched control removes only the energy gradient.
These are new-model settings, not a claim to reproduce the historical
architecture or training objective.

Seed 20261112 is primary; 20261113 and 20261114 measure sensitivity.
The protocol, cohorts, development records and fitted models are frozen
before independent evaluation. No test outcome informs thresholds,
checkpoints or seed choice.
